\documentclass[12pt, titlepage]{article}

\usepackage{booktabs}
\usepackage{tabularx}
\usepackage{hyperref}
\usepackage{float}
\usepackage{placeins}
\usepackage{longtable}
\usepackage{caption}
\usepackage{array}

\hypersetup{
    colorlinks,
    citecolor=blue,
    filecolor=black,
    linkcolor=red,
    urlcolor=blue
}
\usepackage[round]{natbib}

\input{../Comments}
%% Common Parts

\newcommand{\progname}{Project Proxi} % PUT YOUR PROGRAM NAME HERE
\newcommand{\authname}{Team \#23, Project Proxi
\\ Savinay Chhabra
\\ Amanbeer Singh Minhas
\\ Gourob Podder
\\ Ajay Singh Grewal} % AUTHOR NAMES                  

\usepackage{hyperref}
    \hypersetup{colorlinks=true, linkcolor=blue, citecolor=blue, filecolor=blue,
                urlcolor=blue, unicode=false}
    \urlstyle{same}
                                


\begin{document}

\title{System Verification and Validation Plan for \progname{}}
\author{\authname}
\date{\today}

\maketitle

\pagenumbering{roman}

\section*{Revision History}

\begin{tabularx}{\textwidth}{p{2.3cm}p{1.1cm}X}
\toprule
\textbf{Date} & \textbf{Ver.} & \textbf{Notes} \\
\midrule
30/10/2025 & 1.0 & Added initial V\&V plan. \\
08/03/2026 & 1.1 & Added introductory text for Section 1. \\
 &  & Corrected spelling, grammar, and document flow issues. \\
 &  & Replaced magic numbers with symbolic constants where appropriate. \\
 &  & Standardized terminology and stylization (for example, macOS). \\
 &  & Moved table titles above tables and added in-text references. \\
 &  & Added a role reassignment plan for V\&V team members. \\
 &  & Refined review checklists so items are atomic and directly checkable. \\
 &  & Clarified specific WCAG checks used for accessibility verification. \\
 &  & Added tests for unintended and unauthorized actions, including \\
 &  & destructive and replay-style scenarios. \\
 &  & Included draft checklists referenced by the plan. \\
 &  & Clarified traceability with requirement-to-test matrix coverage. \\
 &  & Made walkthroughs, inspections, stakeholder reviews, and \\
 &  & project-specific extras more explicit and justified. \\
\bottomrule
\end{tabularx}

\newpage

\tableofcontents

\newpage

\pagenumbering{arabic}

\noindent This document outlines the Verification and Validation plan for
Project Proxi. It explains how our team plans to evaluate the system we are
building against its requirements. The plan describes what tests will be
run, how documents and designs will be reviewed and validated, and how the
evidence collected traces back to the original SRS document. The main goal
is to make sure Project Proxi is reliable, easy to use, and genuinely
helpful for the people who rely on it for accessibility and productivity.

\section{General Information}

\subsection{Summary}
The aim of this document is to present the verification and validation
plan for Project Proxi. Project Proxi is an AI-powered desktop assistant
that allows its users to interact with their computer using voice and
text. The target users for Project Proxi include seniors, individuals
with disabilities, and power users who want a more accessible,
efficient, and convenient way to perform computer tasks.

\subsection{Objectives}
The main objective of the Verification and Validation (V\&V) plan is to
confirm that Project Proxi satisfies its stated functional and
non-functional requirements and fulfills its primary purpose of improving
computer accessibility and efficiency for seniors, individuals with
disabilities, and power users.

\noindent The principal objectives that this V\&V plan aims to meet are:
\begin{itemize}
  \item \textbf{Verify Functional Correctness:} Ensure that the features
  described in the SRS are implemented and behave as expected through
  unit, integration, and system-level testing.
  \item \textbf{Validate Usability and Accessibility:} Demonstrate that
  users within the target demographic can effectively operate the system
  using voice and text while meeting the relevant accessibility
  requirements, including WCAG 2.2 guidance for accessible user
  interfaces.
  \item \textbf{Assess Performance and Reliability:} Confirm that the
  application meets performance, stability, and error-handling goals as
  defined in the SRS document.
  \item \textbf{Hardware Compatibility Testing:} Meet verification goals
  on a variety of consumer hardware devices. Specialized or adaptive
  devices will not be included; only commercially available Windows and
  macOS computers.
\end{itemize}

Out-of-scope objectives include extensive long-term field usability
studies, comprehensive privacy compliance, and specialized hardware
testing. The project will assume that any external libraries have
already been verified by their respective development teams. These
exclusions are justified given the project's limited time, budget, and
academic context. Verification efforts are therefore directed toward
objectives that help meet core software requirements.

\subsection{Challenge Level and Extras}

This project does not include a designated challenge level, as its
primary focus is on the development and validation of a functional,
accessible, and user-friendly AI-powered desktop assistant. The project
scope aligns with the general expectations of the capstone course and
emphasizes reliability, usability, and compliance with accessibility
standards rather than advanced AI research or complex algorithmic
innovation.

\noindent This project includes the following approved extras:
\begin{itemize}
  \item \textbf{User Manual:} A user manual will document installation,
  setup, and everyday usage of the software through clear, step-by-step
  instructions. This extra is important because Project Proxi is intended
  for seniors, individuals with disabilities, and power users with
  varying levels of technical proficiency. The manual supports onboarding
  and independent use by explaining core workflows, common commands,
  safety confirmations, and troubleshooting guidance in a structured and
  easy-to-reference format.
  \item \textbf{Usability Report:} A usability report will be developed
  to assess the system's core objectives and requirements. It will
  summarize the results of testing as specified in the SRS document.
  This extra is important because usability and accessibility are
  central goals of Project Proxi, and the report will help identify
  issues that can be used to improve the system before final delivery.
\end{itemize}

\subsection{Relevant Documentation}

The following documentation is directly relevant to the verification and
validation activities for Project Proxi:
\begin{itemize}
  \item \textbf{Development Plan [\citet{DevPlan}]:} This document
  defines the overall project workflow, timelines, and milestones for
  the different phases of the project.
  \item \textbf{Problem Statement and Goals [\citet{ProbStatement}]:}
  This document defines the problems, goals, and stakeholders of the
  project.
  \item \textbf{Software Requirements Specification [\citet{SRS}]:}
  Describes the functional and non-functional requirements of the
  project.
  \item \textbf{Hazard Analysis [\citet{HazardAnalysis}]:} This
  document identifies any potential sources of harm or risk that may be
  associated with the project. Potential risks such as unauthorized
  access, misuse, or data breaches and their mitigation strategies are
  documented here.
  \item \textbf{Usability Report [\citet{UsabilityReport}]:} Includes
  results of user evaluation and beta testing during the validation
  stage. It summarizes key findings related to the system meeting its
  design and usability goals.
  \item \textbf{User Manual [\citet{UserManual}]:} Documents how to
  install, configure, and use the software, including command examples,
  core workflows, and troubleshooting steps to support users with
  varying technical proficiency.
\end{itemize}

\section{Plan}
This section outlines the strategy for Verification and Validation (V\&V)
across all phases of the Project Proxi lifecycle. The plan defines the roles
and responsibilities of the V\&V team and specifies structured approaches for
verifying the Software Requirements Specification (SRS), Design, and
Implementation, as well as the approach for Software Validation. This
systematic process increases confidence in the correctness, accessibility,
and reliability of the final system.

Verification ensures that the system is built correctly and meets its stated
requirements, while validation confirms that the right system has been built
for its intended users. Both static techniques (reviews, inspections, and
walkthroughs) and dynamic techniques (testing and user evaluations) are used.
All artifacts, results, and issues are tracked through GitHub for full
traceability between requirements, test cases, and evidence.

\subsection{Verification and Validation Team}

The Verification and Validation team for Project Proxi is the Proxi project
team itself. Since the team is small, all members contribute to V\&V
activities in addition to their development responsibilities. Table~
\ref{tab:vnv-team} summarizes the primary V\&V responsibilities assigned to
each team member.

All verification tasks and results are peer reviewed, documented, and
tracked in GitHub to maintain accountability and consistent quality
assurance.

\begin{table}[H]
\caption{Verification and Validation Team Roles}
\label{tab:vnv-team}
\centering
\begin{tabular}{|p{3cm}|p{3.5cm}|p{7.5cm}|}
\hline
\textbf{Name} & \textbf{Role} & \textbf{V\&V Responsibilities} \\
\hline
Savinay Chhabra & Team Lead / Project Manager &
Coordinates all V\&V activities, schedules document and design reviews,
ensures traceability matrices are complete, and validates plan
feasibility. \\
\hline
Amanbeer Singh Minhas & Verification Engineer / Documentation Lead &
Authors test specifications, manages accessibility validation
(WCAG~2.2), executes functional and non-functional tests, and maintains
V\&V evidence and metrics. \\
\hline
Gourob Podder & AI / Backend Developer &
Verifies the accuracy and latency of voice, text, and natural language
processing modules, conducts integration and performance tests, and
performs code inspections. \\
\hline
Ajay Singh Grewal & Front-End / Integration Engineer &
Implements automated UI tests, validates responsiveness and error
handling, and ensures cross-platform consistency and workflow
reliability. \\
\hline
\end{tabular}
\end{table}

\subsection{Role Reassignment and Contingency Plan}

If timelines, availability, or project risks change, V\&V
responsibilities will be rebalanced using the following procedure:
\begin{itemize}
  \item The Team Lead identifies blocked work items and proposes
  reassignment within one business day.
  \item At least one backup reviewer is maintained for each critical
  activity, including SRS review, design review, test execution, and
  evidence tracking.
  \item Reassignments are recorded in a GitHub issue labeled
  \texttt{vnv-role-change}, including the new owner, due date, and
  rationale.
  \item The traceability matrix and meeting notes are updated in the
  same sprint to reflect responsibility changes.
\end{itemize}

\subsection{SRS Verification}

The Software Requirements Specification (SRS) will be verified using a
structured review process involving internal checks, peer feedback, and
consistency checks. The goal is to confirm that all requirements are clear,
consistent, complete, and ready for design and testing.

\vspace{0.5em}
\noindent\textbf{Planned Verification Stages:}
\begin{itemize}
\item \textbf{Internal Team Review:}
  Each team member reviews assigned SRS sections to ensure requirements are
  clearly written, uniquely numbered, feasible, and testable. Feedback and
  questions are added as GitHub issues under the \texttt{SRS-review} tag for
  tracking.

\item \textbf{Peer Team Feedback:}
  Another project team will review the SRS to identify unclear, ambiguous,
  inconsistent, or missing requirements. Their comments will be discussed in
  a short review meeting, and agreed changes will be included in the next SRS
  revision.

\item \textbf{Consistency and Traceability Check:}
  Manual and automated checks will confirm that every requirement label
  (functional, non-functional, and constraint-related) appears in the
  traceability matrix and links to one or more planned verification
  activities in the V\&V Plan. Any missing, duplicate, or inconsistent
  identifiers will be flagged for correction.

\item \textbf{Team Validation Meeting:}
  After all updates, the team will meet to confirm that project goals such as
  accessibility, reliability, safety, and usability are fully captured in the
  SRS. The final version will then be approved as the baseline for design and
  testing.
\end{itemize}

\vspace{0.5em}
\noindent\textbf{SRS Review Checklist:}
\begin{itemize}
\item Each requirement has an identifier.
\item Each identifier is unique.
\item Each requirement contains one main verifiable idea.
\item Each requirement uses direct language.
\item Undefined terms are avoided or added to the glossary.
\item Each requirement states observable behaviour or a measurable outcome.
\item Each requirement is feasible within capstone scope.
\item Each requirement can be verified by inspection, analysis, or test.
\item Each requirement is consistent with the problem statement.
\item Each requirement is consistent with the project goals.
\item Each safety-related requirement is consistent with the hazard analysis.
\item Accessibility requirements reference WCAG~2.2 where applicable.
\item Voice interaction requirements are represented where needed.
\item Text interaction requirements are represented where needed.
\item High-risk actions include confirmation or mitigation requirements.
\item Each requirement appears in the traceability matrix.
\end{itemize}

\subsection{Design Verification}

The design of Project Proxi will be verified through focused peer reviews,
walkthroughs, and checklists to confirm that the architecture, interfaces,
and data flows align with the approved SRS. Each review will help identify
inconsistencies or missing details before implementation begins.

\vspace{0.5em}
\noindent\textbf{Planned Verification Activities:}
\begin{itemize}
\item \textbf{Internal Design Review:}
  Team members will examine assigned modules to ensure the design is modular,
  consistent with SRS requirements, and feasible for implementation. Each
  reviewer will complete a checklist covering module responsibilities,
  interface definitions, traceability, and error handling.

\item \textbf{Cross-Team Design Review:}
  Another project team will review the design documents and provide feedback
  on readability, logic, completeness, and consistency. Their comments will
  be logged as GitHub issues tagged \texttt{Design-review} and resolved in
  the next revision.

\item \textbf{Walkthrough Meeting:}
  The team will conduct a collaborative walkthrough in which each member
  explains the purpose, inputs, outputs, and dependencies of their assigned
  subsystem. Meeting notes will record identified issues, action items, and
  owners.

\item \textbf{Traceability Check:}
  A traceability matrix will confirm that every major design component maps to
  at least one SRS requirement and that every SRS requirement is supported by
  one or more design elements. Any unmatched modules or requirements will be
  flagged for revision before coding starts.
\end{itemize}

\vspace{0.5em}
\noindent\textbf{Design Review Checklist Items:}
\begin{itemize}
\item Each module has a unique name.
\item Each module states its main purpose.
\item Each module lists its responsibilities.
\item Each module lists its inputs.
\item Each module lists its outputs.
\item Each interface names the caller module.
\item Each interface names the callee module.
\item Each interface lists exchanged data items.
\item Data item names match across connected modules.
\item Data item meaning is consistent across connected modules.
\item Each major data flow identifies its source.
\item Each major data flow identifies its destination.
\item Each major data flow includes the expected transformation or decision.
\item Error handling is identified where failure is possible.
\item Overlapping module responsibilities are avoided.
\item Accessibility-related design elements are identified explicitly.
\item Voice interaction support appears in the design.
\item Text interaction support appears in the design.
\item Each major component traces to at least one SRS requirement.
\end{itemize}

\subsection{Verification and Validation Plan Verification}

The Verification and Validation (V\&V) Plan itself will also be verified to
ensure it is complete, consistent, and realistic. This review confirms that
the plan defines clear responsibilities, test coverage, and measurable
evidence of software quality.

\vspace{0.5em}
\noindent\textbf{Planned Verification Activities:}
\begin{itemize}
\item \textbf{Internal Review:}
  Team members will inspect the V\&V Plan for completeness, logical flow, and
  consistency with other project documents such as the SRS, MIS, MG, and
  Hazard Analysis. Issues will be logged on GitHub under the tag
  \texttt{VnV-review}.

\item \textbf{Peer Team Review:}
  A partner project team will review the plan to provide external feedback on
  clarity, organization, feasibility, and coverage. Comments will be
  discussed in a team meeting, and accepted suggestions will be incorporated
  into the next version.

\item \textbf{Cross-Document Consistency Check:}
  The team will confirm that all V\&V activities listed here are consistent
  with the SRS, MIS, MG, and Hazard Analysis and that requirement identifiers
  match across artifacts.

\item \textbf{Review Process Validation:}
  Small intentional edits, such as removing a trace link or altering a test
  identifier, will be introduced to confirm that the team's review process
  can detect inconsistencies. This activity validates the effectiveness of
  the review workflow.
\end{itemize}

\vspace{0.5em}
\noindent\textbf{V\&V Plan Review Checklist:}
\begin{itemize}
\item Each cited document title is correct.
\item Each cited document identifier is correct.
\item Roles and responsibilities match Table~\ref{tab:vnv-team}.
\item Each major section includes a short purpose statement.
\item Each planned activity identifies expected evidence.
\item Each test identifier is unique.
\item Each test identifier is used consistently across the document.
\item Internal review steps are stated explicitly.
\item Peer review steps are stated explicitly.
\item The plan fits the project timeline and tool availability.
\end{itemize}

\subsection{Implementation Verification}

The implementation of Project Proxi will be verified through a combination of
dynamic and static techniques to ensure the system performs correctly,
reliably, and according to its specified design. This section summarizes how
unit, integration, and system testing activities will be supported by
automated tools and structured reviews.

\vspace{0.5em}
\noindent\textbf{Dynamic Verification:}
\begin{itemize}
\item \textbf{Unit Testing:}
  Each software module will have dedicated unit tests as outlined in the unit
  testing plan. Tests will confirm correct outputs, data flow, state changes,
  and error handling for input processing, planning, execution, and UI
  modules. Coverage results will be monitored through automated reports to
  ensure that major functions are exercised.

\item \textbf{Integration and System Testing:}
  Combined module testing will validate interactions between the input
  pipeline, user interface, planning components, execution components, and
  backend services. These tests correspond directly to the functional and
  non-functional test cases listed in this document.

\item \textbf{Mutation Testing:}
  To evaluate the quality of the test suite, small artificial faults
  (mutants) will be introduced into the codebase. Existing unit and system
  tests will be rerun to confirm that they detect the injected faults. A high
  mutation score will increase confidence in the strength and completeness of
  the test suite.
\end{itemize}

\vspace{0.5em}
\noindent\textbf{Static Verification:}
\begin{itemize}
\item \textbf{Code Walkthroughs and Inspections:}
  Code walkthroughs will occur before major feature merges or at least once
  per sprint for active modules. The author will present the change, linked
  design elements, linked test cases, and expected behaviour. Review notes,
  blocking issues, and follow-up actions will be recorded in pull request
  comments or GitHub issues. A walkthrough is considered complete only when
  all blocking issues have an owner and a resolution plan.

\item \textbf{Static Analysis Tools:}
  Linters and analyzers, such as ESLint, Pylint, or similar tools, will be
  used to detect common issues, unused variables, inconsistent style, and
  potential logic errors. Results will be documented as part of the
  verification evidence.

\item \textbf{Code Review via Pull Requests:}
  All implementation changes will go through GitHub pull requests with at
  least one reviewer. Reviewers will verify coding standards, style,
  alignment with design diagrams, and whether the change preserves intended
  behaviour without introducing unintended actions.
\end{itemize}

\subsection{Automated Testing and Verification Tools}

Automated tools will be used throughout the implementation phase to support
testing, continuous integration, and quality assurance. These tools help
reduce manual effort, detect regressions early, and provide measurable
evidence of software reliability and maintainability.

\vspace{0.5em}
\noindent\textbf{Automated Testing Framework:}
\begin{itemize}
\item The project will use \texttt{pytest} for automated unit,
  integration, and scripted scenario tests. Each major module will have
  dedicated test files and fixtures to verify inputs, outputs, state
  changes, and edge cases.
\item Scripted scenario tests will simulate realistic user interactions by
  submitting selected voice and text commands, capturing system responses,
  and comparing them against expected results.
\item For repeatable testing, scenario tests will use a predefined command
  set maintained by the team. The full set of commands used for automated
  testing will be listed explicitly in the corresponding test cases or in an
  appendix, rather than described only as a sample count.
\end{itemize}

\vspace{0.5em}
\noindent\textbf{Continuous Integration and Delivery (CI/CD):}
\begin{itemize}
\item GitHub Actions will be configured to automatically run linting and
  automated tests whenever code is pushed or a pull request is opened.
  The pipeline will help detect regressions early and ensure that changes
  are reviewed against a consistent verification process before merge.
\end{itemize}

\vspace{0.5em}
\noindent\textbf{Static Analysis and Linters:}
\begin{itemize}
\item Static analysis tools such as \texttt{pylint} and \texttt{flake8}
  will be used to detect syntax errors, unused variables, inconsistent
  style, and maintainability issues. Results will be logged and tracked
  as part of the verification evidence.
\end{itemize}

\vspace{0.5em}
\noindent\textbf{Build and Automation Utilities:}
\begin{itemize}
\item The project will use \texttt{make} commands and supporting scripts
  to automate environment setup, dependency installation, linting, and
  test execution. This helps ensure consistent verification results
  across development machines and the CI/CD pipeline.
\end{itemize}

\subsection{Software Validation}

Software validation confirms that Project Proxi satisfies user needs and
performs as intended in realistic scenarios. The focus is on showing that
the system achieves its goals of accessibility, responsiveness, safety,
and ease of use for seniors, individuals with disabilities, and power
users.

\vspace{0.5em}
\noindent\textbf{Planned Validation Activities:}
\begin{itemize}
\item \textbf{User Testing Sessions:}
  A small group of peers and volunteers will perform task-based
  evaluations, such as launching applications, entering text, and
  managing files through voice and text interaction. Their feedback will
  be used to assess ease of learning, efficiency, accessibility, and
  overall satisfaction.

\item \textbf{Accessibility Validation:}
  Validation will confirm that the interface supports the accessibility
  goals defined in the SRS and aligns with relevant WCAG~2.2 guidance.
  The evaluation will explicitly check WCAG~2.2 Success Criteria 1.4.3
  (Contrast (Minimum)), 1.4.4 (Resize Text), 1.3.1 (Info and
  Relationships), and 4.1.2 (Name, Role, Value) where applicable to the
  final interface.

\item \textbf{Stakeholder Review Session:}
  The team will conduct one structured review session as a validation
  milestone. The participants will include course staff, classmates, and
  peer reviewers acting as informed external stakeholders. The session
  will include a short system overview, a guided demonstration of core
  tasks, a demonstration of one high-risk confirmation flow, and a brief
  feedback period. The core demonstration scenarios will include email
  summarization, calendar event creation, browser search and summary,
  Notion page creation, and a privileged action that requires explicit
  confirmation. Feedback will be collected through a short review form
  and meeting notes. Any issue reported by at least two reviewers, or
  any issue that blocks a core task, will be logged as a GitHub issue
  and considered for revision before final delivery.

\item \textbf{Automated Scenario Validation:}
  Scripted scenario tests will submit a predefined and fully enumerated
  set of commands to the system, record the observed outputs, and compare
  them with expected behaviour. These tests will help validate recurring
  tasks, command interpretation accuracy, response consistency, and
  protection against unintended actions.

\item \textbf{External Data and Comparison:}
  Where appropriate, publicly available speech or language datasets may
  be used to compare command interpretation behaviour against baseline
  expectations. This provides an additional reference point for judging
  consistency and robustness.
\end{itemize}

\section{System Tests}

Each test is specific and measurable, providing sufficient detail for
replication once the system is built. The tests are grouped into logical
areas covering input handling, natural language understanding, task
execution, feedback mechanisms, memory management, logging, and system
safety.

Unless otherwise stated, repeated thresholds and sample sizes reference the
symbolic parameters defined in Appendix~A. These thresholds are chosen to
align with SRS fit criteria and the realistic limits of a capstone project,
rather than assuming perfect behaviour from speech- and language-based
features.

\subsection{Tests for Functional Requirements}

The following tests verify that each functional requirement defined in the
SRS has been properly implemented and behaves as intended. Each test
specifies the type of testing, initial system state, inputs, expected
results, and how the test will be performed.

\subsubsection{Input and Recognition Requirements}

\begin{enumerate}

\item \textbf{Test FUNC.R.1 -- Speech and Text Input Acceptance}

\textbf{Type:} Functional / Manual Testing

\textbf{Initial State:} System initialized and idle at the main
interface.

\textbf{Input/Condition:} Users provide the following
\texttt{CMD\_SET\_SIZE} valid commands. Commands 1--10 are entered by
text, and commands 11--20 are spoken through the microphone interface.

{\renewcommand{\labelenumii}{\arabic{enumii}.}
\begin{enumerate}
  \item Summarize my latest unread email.
  \item Draft an email to John about tomorrow's team meeting.
  \item Show my calendar events for tomorrow.
  \item Create a calendar event for Friday at 2 PM called Team Sync.
  \item Search the web for Hamilton weather and summarize the result.
  \item Search the web for nearby walk-in clinics.
  \item Make a Notion page called Drinking Water Tracker.
  \item Add ``Drink 8 glasses today'' to my Drinking Water Tracker page.
  \item Summarize the file named Capstone Notes.
  \item What can I say?
  \item Summarize my latest unread email.
  \item Draft an email to John about tomorrow's team meeting.
  \item Show my calendar events for tomorrow.
  \item Create a calendar event for Friday at 2 PM called Team Sync.
  \item Search the web for Hamilton weather and summarize the result.
  \item Search the web for nearby walk-in clinics.
  \item Make a Notion page called Drinking Water Tracker.
  \item Add ``Drink 8 glasses today'' to my Drinking Water Tracker page.
  \item Summarize the file named Capstone Notes.
  \item What can I say?
\end{enumerate}
}

\textbf{Output/Result:} The system correctly accepts and interprets at
least \texttt{CMD\_ACCEPT\_MIN} of the valid inputs, and each accepted
input is converted into the expected internal command representation and
displayed result.

\textbf{How test will be performed:} Each input and resulting interpreted
command will be logged automatically. The test supervisor will compare
the interpreted result against the expected result for the full command
set and calculate the success rate. Any rejected, misinterpreted, or
malformed results will be recorded for analysis.

\textbf{Functional Requirements Covered:} FUNC.R.1

\item \textbf{Test FUNC.R.2 -- Speech-to-Text Accuracy}

\textbf{Type:} Functional / Automated Evaluation

\textbf{Initial State:} System with speech-to-text enabled and
configured in a controlled acoustic environment.

\textbf{Input/Condition:} The following \texttt{STT\_CMD\_COUNT} spoken
commands are each recorded from \texttt{STT\_SPEAKER\_COUNT} speakers:

{\renewcommand{\labelenumii}{\arabic{enumii}.}
\begin{enumerate}
  \item Summarize my latest unread email.
  \item Draft an email to John about tomorrow's team meeting.
  \item Show my calendar events for tomorrow.
  \item Create a calendar event for Friday at 2 PM called Team Sync.
  \item Search the web for Hamilton weather and summarize the result.
  \item Search the web for nearby walk-in clinics.
  \item Make a Notion page called Drinking Water Tracker.
  \item Add ``Drink 8 glasses today'' to my Drinking Water Tracker page.
  \item Summarize the file named Capstone Notes.
  \item What can I say?
\end{enumerate}
}

\textbf{Output/Result:} The resulting transcriptions achieve at least
\texttt{STT\_ACC\_MIN} word accuracy across the full sample set.

\textbf{How test will be performed:} The recorded audio samples will be
batch fed into the speech-to-text component. Output transcripts will be
compared against reference transcripts using a word error rate metric,
and the overall accuracy will be computed from the full set.

\textbf{Functional Requirements Covered:} FUNC.R.2

\end{enumerate}

\subsubsection{Intent Understanding and Task Execution}

\begin{enumerate}

\item \textbf{Test FUNC.R.3 -- Intent Recognition Accuracy}

\textbf{Type:} Functional / Automated Testing

\textbf{Initial State:} Natural language understanding component
initialized with the current intent parsing configuration.

\textbf{Input/Condition:} The system is given
\texttt{INTENT\_SAMPLE\_COUNT} labeled inputs covering the supported
intents. The supported intents are: summarize email, draft email, show
calendar, create calendar event, browser search, summarize browser
result, create Notion page, update Notion page, summarize file, and show
help or examples.

\textbf{Output/Result:} At least \texttt{INTENT\_ACC\_MIN} of the
labeled inputs are assigned the correct ground-truth intent.

\textbf{How test will be performed:} A test script will feed the labeled
inputs into the intent parser. Predicted intents will be logged and
compared against human-labeled references to compute intent-level
accuracy.

\textbf{Functional Requirements Covered:} FUNC.R.3

\item \textbf{Test FUNC.R.4 -- Task Planning and Execution Success Rate}

\textbf{Type:} Functional / Integration Testing

\textbf{Initial State:} System active with required tools, permissions,
and network connectivity available.

\textbf{Input/Condition:} The following representative high-level
commands are executed sequentially:

{\renewcommand{\labelenumii}{\arabic{enumii}.}
\begin{enumerate}
  \item Summarize my latest unread email.
  \item Draft an email to Sarah about the capstone update.
  \item Draft an email to John saying I will join in 10 minutes.
  \item Show my calendar events for tomorrow.
  \item Create a calendar event for Friday at 2 PM called Team Sync.
  \item Move my Team Sync event to 3 PM.
  \item Search the web for Hamilton weather.
  \item Summarize the current browser result.
  \item Search the web for nearby pharmacies.
  \item Summarize the top result in two sentences.
  \item Make a Notion page called Drinking Water Tracker.
  \item Add ``Drink 8 glasses today'' to my Drinking Water Tracker page.
  \item Make a Notion page called Capstone Tasks.
  \item Add ``Finish V\&V plan'' to my Capstone Tasks page.
  \item Summarize the file named Capstone Notes.
  \item Read the first paragraph of Capstone Notes aloud.
  \item Show the summary again.
  \item What can I say?
  \item Undo the last action.
  \item Cancel the current action.
\end{enumerate}
}

\textbf{Output/Result:} At least \texttt{TASK\_SUCCESS\_MIN} of the
commands are completed successfully and produce the intended task outcome
in the interface.

\textbf{How test will be performed:} Commands will be executed one at a
time. Testers will verify whether each displayed result, draft, summary,
or state change matches the intended user goal. Failures will be logged
along with their cause, such as planning failure, permission failure, or
execution failure.

\textbf{Functional Requirements Covered:} FUNC.R.4

\end{enumerate}

\subsubsection{Feedback and Responsiveness}

\begin{enumerate}

\item \textbf{Test FUNC.R.5 -- Response Feedback Timing}

\textbf{Type:} Functional / Manual Timing

\textbf{Initial State:} System ready to execute user commands.

\textbf{Input/Condition:} The following common tasks are executed:
summarize latest unread email, draft an email, show tomorrow's calendar,
create a calendar event, search the web for Hamilton weather, summarize
the current result, search for nearby pharmacies, make a Notion page,
summarize a file, and show help commands.

\textbf{Output/Result:} Each command generates visible or spoken
feedback within \texttt{FEEDBACK\_DELAY\_MAX} after task completion or
after the system determines that the task cannot be completed.

\textbf{How test will be performed:} A timer will start at the moment the
task finishes or the failure state is known. Feedback appearance time
will be recorded for each trial, and all observed times will be compared
against the threshold.

\textbf{Functional Requirements Covered:} FUNC.R.5

\end{enumerate}

\subsubsection{Contextual Memory and Logging}

\begin{enumerate}

\item \textbf{Test FUNC.R.6 -- Contextual Follow-Up Interpretation}

\textbf{Type:} Functional / Manual Sequence Testing

\textbf{Initial State:} System with active short-term conversational
memory.

\textbf{Input/Condition:} Testers execute the following conversational
sequences:

{\renewcommand{\labelenumii}{\arabic{enumii}.}
\begin{enumerate}
  \item Summarize my latest unread email. Reply to it with thanks.
  \item Draft an email to John. Change the subject to Project Update.
  \item Show my calendar for tomorrow. Add office hours at 4 PM.
  \item Create a calendar event for Friday at 2 PM. Move it to 3 PM.
  \item Search the web for Hamilton weather. Summarize it.
  \item Search the web for nearby walk-in clinics. Open the first result
  and summarize it.
  \item Make a Notion page called Drinking Water Tracker. Add today's
  intake to it.
  \item Make a Notion page called Capstone Tasks. Add testing checklist
  to it.
  \item Summarize the file named Capstone Notes. Read it aloud.
  \item Show my latest drafted email. Send it.
\end{enumerate}
}

\textbf{Output/Result:} At least \texttt{FOLLOWUP\_SUCCESS\_MIN} of the
sequences are completed with correct interpretation of follow-up
references such as ``it'', ``that'', or omitted repeated context.

\textbf{How test will be performed:} Each sequence will be entered by a
tester. The interpreted actions and resulting displayed outputs will be
reviewed to confirm that contextual references are resolved correctly.

\textbf{Functional Requirements Covered:} FUNC.R.6

\item \textbf{Test FUNC.R.7 -- Interaction and Task Logging Verification}

\textbf{Type:} Functional / Inspection

\textbf{Initial State:} Logging module enabled with write access to the
configured storage location.

\textbf{Input/Condition:} The following sessions are performed, with a
mix of successful and failed tasks using both voice and text:
summarize email, draft email, create calendar event, make Notion page,
browser search, file summary, invalid command, send draft then cancel,
browser request with network disabled, and undo last action.

\textbf{Output/Result:} Each session generates a timestamped log entry
containing the received input, interpreted intent, resulting action, and
final outcome.

\textbf{How test will be performed:} After each session, reviewers will
inspect the recorded logs for completeness, correct timestamps, session
identifiers, action summaries, and outcome fields.

\textbf{Functional Requirements Covered:} FUNC.R.7

\end{enumerate}

\subsubsection{Usability and Safety}

\begin{enumerate}

\item \textbf{Test FUNC.R.8 -- High-Level Command Usability}

\textbf{Type:} Usability / Observational Testing

\textbf{Initial State:} Deployed system interface with voice and text
interaction enabled.

\textbf{Input/Condition:} \texttt{USABILITY\_PARTICIPANTS} participants
attempt the following core tasks using Proxi's high-level interface:
summarize an unread email, draft an email, show tomorrow's calendar,
create a calendar event, search the web, summarize a browser result,
make a Notion page, update a Notion page, summarize a file, and use the
help command.

\textbf{Output/Result:} At least \texttt{HIGHLEVEL\_TASK\_MIN} of
attempted tasks are completed successfully without requiring users to
interact with low-level system details such as raw file paths, terminal
commands, or manual system configuration steps.

\textbf{How test will be performed:} Observers will record whether each
task is completed successfully, whether low-level system details were
needed, and where users hesitated or failed. Completion rate will be
computed over the full task set across all participants.

\textbf{Functional Requirements Covered:} FUNC.R.8

\item \textbf{Test FUNC.R.9 -- Privileged Action Confirmation}

\textbf{Type:} Functional / Safety Validation

\textbf{Initial State:} System running with permission-sensitive actions
enabled.

\textbf{Input/Condition:} The following privileged commands are issued by
both voice and text: send the drafted email, delete an email draft,
delete a calendar event, remove a Notion page, and clear saved session
history. For each command, one trial confirms the action and one trial
cancels the action.

\textbf{Output/Result:} Every privileged command triggers a confirmation
prompt before execution. Confirmed actions execute only after explicit
approval, and cancelled actions do not execute.

\textbf{How test will be performed:} Testers will issue each privileged
command twice, once followed by confirmation and once followed by
cancel. Observers will verify that the prompt appears before execution
and that the final result matches the approval or cancellation path.

\textbf{Functional Requirements Covered:} FUNC.R.9

\end{enumerate}

\subsection{Tests for Non-Functional Requirements}

This section defines specific tests used to determine whether Proxi meets
its non-functional requirements. Each test is derived from the SRS and
states how the result will be verified.

\subsubsection{Look and Feel Requirements}

\begin{enumerate}

\item \textbf{Test NF.LF.1 -- UI Conformance and Readability}

\textbf{Type:} Static Inspection / Manual Testing

\textbf{Initial State:} Latest UI build available in the chat-style
frontend.

\textbf{Input/Condition:} Review the main screens, transcript area,
status indicators, and primary controls using the following checklist:
listen control visible, stop or cancel control visible, undo control
visible, status indicator text readable, transcript text readable,
minimal visual clutter in the main path, advanced options hidden from
the default flow, large click targets, and theme toggle present.

\textbf{Output/Result:} All checklist items are present and usable in the
main workflow.

\textbf{How test will be performed:} Two reviewers will inspect the UI
independently, complete the checklist, compare results, and record any
issues with screenshots.

\textbf{Non-Functional Requirements Covered:} APP.1, APP.2, APP.3,
APP.4, APP.5

\item \textbf{Test NF.LF.2 -- Style and Iconography Check}

\textbf{Type:} Static Inspection / Manual Testing

\textbf{Initial State:} Build with spoken output, on-screen text, and
the final icon set enabled.

\textbf{Input/Condition:} Trigger the following outputs and review the
resulting labels, captions, and icons:
welcome message, task success message, task failure message,
confirmation prompt, cancel action, undo action, help response, and one
browser summary response.

\textbf{Output/Result:} Labels are short and plain, spoken output also
appears on screen, and icons are consistent and paired with short text
labels.

\textbf{How test will be performed:} Reviewers will inspect the
displayed strings and visual elements, then record any inconsistencies
with screenshots or screen recordings.

\textbf{Non-Functional Requirements Covered:} STY.1, STY.2, STY.3

\end{enumerate}

\subsubsection{Usability and Humanity Requirements}

\begin{enumerate}

\item \textbf{Test NF.US.1 -- Task-Based Usability and Learnability}
\textbf{Study}

\textbf{Type:} Dynamic / Manual Usability Testing

\textbf{Initial State:} New users receive a
\texttt{ORIENTATION\_TIME} orientation to Proxi.

\textbf{Input/Condition:} Each participant performs the following core
tasks:
summarize an unread email, draft an email, show tomorrow's calendar,
create a calendar event, search the web, summarize a browser result,
make a Notion page, update a Notion page, summarize a file, and use the
``What can I say?'' help feature.

\textbf{Output/Result:} Main actions are discoverable, most everyday
tasks are completed without assistance after the orientation, repeated
tasks become faster, example commands are easy to find, messages are
clear, risky actions are confirmed, and error messages suggest a next
step. Survey targets are interpreted using the acceptance criteria
defined in Appendix~A.

\textbf{How test will be performed:} Observers will measure task
completion, task time, repeated-task improvement, and requests for help.
Participants will also complete the usability survey in Appendix~A. If
a survey target is missed, the team will log the issue, identify the
affected workflow, and revise that part of the interface or system
response before final delivery.

\textbf{Non-Functional Requirements Covered:} EOU.R.1, EOU.R.2,
EOU.R.3, LEA.R.1, LEA.R.2, UAP.R.1, UAP.R.2, UAP.R.3

\item \textbf{Test NF.US.2 -- Accessibility, Personalization, and}
\textbf{Locale}

\textbf{Type:} Dynamic Testing / Accessibility Inspection

\textbf{Initial State:} Build with voice and text enabled; captions,
scalable text, and contrast-compliant themes implemented.

\textbf{Input/Condition:} Execute the following voice-only tasks:
start setup, ask for help, summarize an unread email, show tomorrow's
calendar, create a calendar event, cancel the current action, undo the
last action, and exit the app. Then inspect the interface for the
following accessibility checks: captions shown for spoken responses,
text resizes to \texttt{TEXT\_RESIZE\_REQ} without loss of information,
and color contrast meets WCAG~2.2 Success Criterion 1.4.3. Also verify
that interface structure and controls satisfy WCAG~2.2 Success
Criteria 1.3.1 and 4.1.2 where applicable, and that date, time, and
language outputs follow Canadian English formatting.

\textbf{Output/Result:} All listed tasks can be completed through the
voice-only path, captions are present for spoken responses, text scaling
works without loss of function, contrast is sufficient, interface
semantics remain understandable to assistive technologies, and Canadian
English formats are used consistently.

\textbf{How test will be performed:} Reviewers will run the voice-only
task suite, inspect captions and scaled layouts, check contrast using an
accessibility checker, inspect control labels and roles in the UI, and
verify locale-formatted outputs in the interface.

\textbf{Non-Functional Requirements Covered:} ACC.R.1, ACC.R.2,
PER.R.1, PER.R.2

\end{enumerate}

\subsubsection{Performance and Reliability Requirements}

\begin{enumerate}

\item \textbf{Test NF.PF.1 -- Latency and Recognition Accuracy Study}

\textbf{Type:} Dynamic / Automated Performance Testing

\textbf{Initial State:} Standard hardware environment with typical
system load.

\textbf{Input/Condition:} Run the following commands in a quiet room and
then repeat the same commands in a moderately noisy indoor environment:
summarize an unread email, draft an email, show tomorrow's calendar,
create a calendar event, search the web for Hamilton weather, summarize
the browser result, search for nearby pharmacies, make a Notion page,
summarize a file, and use the help command.

\textbf{Output/Result:} Voice responses begin within the SAL.R.1 limit,
common actions complete within the SAL.R.2 limit, quiet-environment
recognition accuracy meets POA.R.1, and moderate-noise accuracy meets
POA.R.2.

\textbf{How test will be performed:} The team will time each command,
compare recognized inputs against expected inputs, compute command
accuracy for both environments, and record the results in a test log.

\textbf{Non-Functional Requirements Covered:} SAL.R.1, SAL.R.2,
POA.R.1, POA.R.2

\item \textbf{Test NF.PF.2 -- Robustness and Capacity Study}

\textbf{Type:} Dynamic / Automated and Manual Testing

\textbf{Initial State:} System running on standard hardware.

\textbf{Input/Condition:} Submit the following invalid or adverse
inputs:
empty command, incomplete command, unsupported command, malformed text
input, interrupted speech input, network-dependent request with network
disabled, and a cancelled action. Then run the system continuously for
at least \texttt{SESSION\_DURATION\_MIN} under normal usage.

\textbf{Output/Result:} Input errors are logged and reported without a
crash, the system continues running after invalid or incomplete
commands, continuous operation is maintained for at least
\texttt{SESSION\_DURATION\_MIN}, and CPU usage stays within the CAP.R.2
limit during normal operation.

\textbf{How test will be performed:} Reviewers will trigger each listed
failure condition, inspect logs and on-screen errors, monitor whether
the system remains usable, and record runtime and CPU usage during the
continuous run.

\textbf{Non-Functional Requirements Covered:} ROFT.R.1, ROFT.R.2,
CAP.R.1, CAP.R.2

\item \textbf{Test NF.PF.3 -- Safety, Extensibility, and Longevity}
\textbf{Study}

\textbf{Type:} Dynamic / Manual and Checklist-Based Testing

\textbf{Initial State:} System running with undo support, update
support, and MCP tool registry enabled.

\textbf{Input/Condition:} Perform the following actions:
delete a draft and cancel, delete a draft and confirm, remove a
calendar event and undo it, add one new MCP tool through the registry,
update the build, and simulate an update rollback.

\textbf{Output/Result:} Destructive actions require confirmation,
critical actions can be undone in one step, a new feature becomes
available without major changes to existing functionality, user settings
persist after update, and rollback completes within the LON.R.2 limit.

\textbf{How test will be performed:} Reviewers will execute the listed
actions, verify confirmation and undo behaviour, register one new MCP
tool, confirm that it loads correctly, update the system, and then
measure the rollback process.

\textbf{Non-Functional Requirements Covered:} SAF.R.1, SAF.R.2,
SOE.R.1, SOE.R.2, LON.R.1, LON.R.2

\end{enumerate}

\subsubsection{Operational and Environmental Requirements}

\begin{enumerate}

\item \textbf{Test NF.OP.1 -- Environment and Ecosystem Fit}

\textbf{Type:} Dynamic / Checklist-Based Testing

\textbf{Initial State:} Fresh install on current Windows and macOS
systems with required integrations enabled.

\textbf{Input/Condition:} Run the system indoors under normal operating
conditions, including moderate indoor noise, user distance within
\texttt{MIC\_DISTANCE\_MAX} of the microphone, and room temperature
between \texttt{ROOM\_TEMP\_MIN} and \texttt{ROOM\_TEMP\_MAX}. Perform a
browser search, email task, file summary task, and calendar task while
other common desktop applications are also running.

\textbf{Output/Result:} The system operates effectively within the
expected physical environment, speech recognition remains usable within
the specified distance and indoor noise conditions, and automation does
not interfere with other running applications.

\textbf{How test will be performed:} Reviewers will execute the listed
tasks under the specified conditions, note any environmental failures,
and check whether other applications continue to function normally.

\textbf{Non-Functional Requirements Covered:} EPE.R.1, EPE.R.2,
EPE.R.3, EPE.R.4, WER.R.1, WER.R.2

\item \textbf{Test NF.OP.2 -- Interfaces, Packaging, and Release Check}

\textbf{Type:} Dynamic / Inspection and Checklist Testing

\textbf{Initial State:} Release candidate build, release notes, and
deployment package available.

\textbf{Input/Condition:} Perform a clean install, verify external
integrations with email, browser, file, and speech services, inspect
communication protocol use, inspect external interaction logs, and
review the packaged release.

\textbf{Output/Result:} Interfaces with adjacent systems use the
expected protocols and logging, the product installs with minimal setup,
the release follows the versioning scheme, and the release package is
version controlled and archived.

\textbf{How test will be performed:} Reviewers will perform a clean
install, execute one task per integration, inspect logs for external
interactions, and review release artifacts, tags, and packaged files.

\textbf{Non-Functional Requirements Covered:} IAS.R.1, IAS.R.2,
IAS.R.3, IAS.R.4, PRD.R.1, PRD.R.2, PRD.R.3, REL.R.1, REL.R.2

\end{enumerate}

\subsubsection{Maintainability and Support Requirements}

\begin{enumerate}

\item \textbf{Test NF.MS.1 -- Onboarding, Support, and Adaptability}

\textbf{Type:} Dynamic / Checklist-Based Testing

\textbf{Initial State:} Fresh clone of the repository, README present,
Report Problem feature present, and configuration file editable.

\textbf{Input/Condition:} A new contributor follows the README to build
and run the project, adds one new MCP tool using a module and registry
entry, generates a support bundle through Report Problem, and changes
the default app configuration without changing code.

\textbf{Output/Result:} The build and run process succeeds without help,
a new MCP tool is added without modifying existing tools, the support
bundle contains recent logs and basic system information with consent,
and the system continues working after only a configuration change.

\textbf{How test will be performed:} The team will time the setup,
verify that the new contributor can complete the README steps, add and
register one tool, inspect the generated ZIP bundle, modify the default
app configuration, and rerun the affected task flow.

\textbf{Non-Functional Requirements Covered:} MTN.R.1, MTN.R.2,
SUP.R.1, SUP.R.2, ADP.R.1, ADP.R.2

\end{enumerate}

\subsubsection{Security Requirements}

\begin{enumerate}

\item \textbf{Test NF.SEC.1 -- Access, Integrity, Privacy, and}
\textbf{Immunity}

\textbf{Type:} Dynamic / Inspection and Security Scanning

\textbf{Initial State:} Fresh install, privacy policy available, and
dependency scanner configured.

\textbf{Input/Condition:} Verify that no account is required to use the
app, trigger actions that require explicit permission, inspect stored
and transmitted data for protection and minimization, review user
consent steps for third-party services, and run dependency
vulnerability scans.

\textbf{Output/Result:} No registration is required, required
permissions must be explicitly granted by the user, stored or
transmitted data is protected from unauthorized alteration, stored data
handling aligns with PIPEDA expectations, third-party data sharing is
minimized and consented to, and no known dependency vulnerabilities
remain unresolved at release time.

\textbf{How test will be performed:} Reviewers will execute permission
flows, inspect privacy-related behaviour and consent prompts, inspect
logs and network traffic at a high level, and run software composition
and dependency scans.

\textbf{Non-Functional Requirements Covered:} ACS.R.1, ACS.R.2, INT.R.1,
PRIV.R.1, IMM.R.1, IMM.R.2

\item \textbf{Test NF.SEC.2 -- Replay or Unauthorized Voice Command}

\textbf{Type:} Dynamic / Manual Security Testing

\textbf{Initial State:} System running with voice input enabled.

\textbf{Input/Condition:} Play a recorded high-risk voice command, or
have a different person issue a destructive command such as deleting a
draft or removing an event.

\textbf{Output/Result:} The system does not execute the action silently.
It must either reject the command, require explicit confirmation, or log
the event for review.

\textbf{How test will be performed:} Reviewers will issue a risky
command using a recording and then using a second speaker, and will
verify that the action is not completed without the expected safeguard.

\textbf{Non-Functional Requirements Covered:} SAF.R.1, ACS.R.2

\end{enumerate}

\subsubsection{Cultural Requirements}

\begin{enumerate}

\item \textbf{Test NF.CUL.1 -- Language and Tone Check}

\textbf{Type:} Static Inspection / Manual Testing

\textbf{Initial State:} Final strings, prompts, and language settings
available.

\textbf{Input/Condition:} Review common screens and messages, including
welcome text, help text, error messages, confirmation prompts, task
success messages, and cancellation messages. Then switch languages and
repeat a subset of basic flows.

\textbf{Output/Result:} Wording remains polite and culturally neutral,
avoids slang and confusing references, and multiple languages support
the basic interaction flow.

\textbf{How test will be performed:} Two reviewers will inspect the
default language strings, then one reviewer will switch languages,
repeat the selected flows, and record any broken or untranslated text.

\textbf{Non-Functional Requirements Covered:} CULR.R.1, CULR.R.2

\end{enumerate}

\subsubsection{Compliance Requirements}

\begin{enumerate}

\item \textbf{Test NF.COM.1 -- Legal and Standards Compliance Check}

\textbf{Type:} Static Inspection / Checklist Testing

\textbf{Initial State:} License file, privacy policy, accessibility
notes, and release artifacts available in the repository.

\textbf{Input/Condition:} Review legal and standards compliance against
the following checklist:
PIPEDA handling for stored user data, MIT license presence and wording,
WCAG~2.2 accessibility evidence, and documentation of software quality
and security practices aligned with ISO~25010 and ISO~27001 goals.

\textbf{Output/Result:} Required legal notices are present, the release
uses the MIT license correctly, privacy documentation aligns with stated
behaviour, and standards-compliance evidence is present and consistent
with the implementation.

\textbf{How test will be performed:} Reviewers will inspect policies,
license files, release notes, and accessibility evidence, then record
any missing or inconsistent items in a compliance checklist.

\textbf{Non-Functional Requirements Covered:} LGL.R.1, LGL.R.2,
STDCOMP.R.1, STDCOMP.R.2

\end{enumerate}

\subsection{Traceability Between Test Cases and Requirements}
\label{section:5.3}

Tables~\ref{tab:fr-trace}, \ref{tab:nfr-trace-a}, and
\ref{tab:nfr-trace-b} show that every requirement defined in the SRS
is covered by at least one planned system test.

\begin{table}[H]
\centering
\small
\caption{Functional Requirements Traceability Matrix}
\label{tab:fr-trace}
\begin{tabularx}{\textwidth}{|c|*{9}{>{\centering\arraybackslash}X|}}
\hline
Test ID & F1 & F2 & F3 & F4 & F5 & F6 & F7 & F8 & F9 \\
\hline
FUNC.R.1 & $\times$ & & & & & & & & \\
\hline
FUNC.R.2 & & $\times$ & & & & & & & \\
\hline
FUNC.R.3 & & & $\times$ & & & & & & \\
\hline
FUNC.R.4 & & & & $\times$ & & & & & \\
\hline
FUNC.R.5 & & & & & $\times$ & & & & \\
\hline
FUNC.R.6 & & & & & & $\times$ & & & \\
\hline
FUNC.R.7 & & & & & & & $\times$ & & \\
\hline
FUNC.R.8 & & & & & & & & $\times$ & \\
\hline
FUNC.R.9 & & & & & & & & & $\times$ \\
\hline
\end{tabularx}
\end{table}
\FloatBarrier

\noindent
Functional requirement keys:
F1 = FUNC.R.1, F2 = FUNC.R.2, F3 = FUNC.R.3, F4 = FUNC.R.4,
F5 = FUNC.R.5, F6 = FUNC.R.6, F7 = FUNC.R.7, F8 = FUNC.R.8,
F9 = FUNC.R.9.

\begin{table}[H]
\centering
\small
\caption{Non-Functional Requirements Traceability Matrix Part 1}
\label{tab:nfr-trace-a}
\begin{tabularx}{\textwidth}{|c|*{6}{>{\centering\arraybackslash}X|}}
\hline
Req ID & LF1 & LF2 & US1 & US2 & PF1 & PF2 \\
\hline
APP.1 & $\times$ & & & & & \\
\hline
APP.2 & $\times$ & & & & & \\
\hline
APP.3 & $\times$ & & & & & \\
\hline
APP.4 & $\times$ & & & & & \\
\hline
APP.5 & $\times$ & & & & & \\
\hline
STY.1 & & $\times$ & & & & \\
\hline
STY.2 & & $\times$ & & & & \\
\hline
STY.3 & & $\times$ & & & & \\
\hline
EOU.R.1 & & & $\times$ & & & \\
\hline
EOU.R.2 & & & $\times$ & & & \\
\hline
EOU.R.3 & & & $\times$ & & & \\
\hline
LEA.R.1 & & & $\times$ & & & \\
\hline
LEA.R.2 & & & $\times$ & & & \\
\hline
UAP.R.1 & & & $\times$ & & & \\
\hline
UAP.R.2 & & & $\times$ & & & \\
\hline
UAP.R.3 & & & $\times$ & & & \\
\hline
ACC.R.1 & & & & $\times$ & & \\
\hline
ACC.R.2 & & & & $\times$ & & \\
\hline
PER.R.1 & & & & $\times$ & & \\
\hline
PER.R.2 & & & & $\times$ & & \\
\hline
SAL.R.1 & & & & & $\times$ & \\
\hline
SAL.R.2 & & & & & $\times$ & \\
\hline
POA.R.1 & & & & & $\times$ & \\
\hline
POA.R.2 & & & & & $\times$ & \\
\hline
ROFT.R.1 & & & & & & $\times$ \\
\hline
ROFT.R.2 & & & & & & $\times$ \\
\hline
CAP.R.1 & & & & & & $\times$ \\
\hline
CAP.R.2 & & & & & & $\times$ \\
\hline
\end{tabularx}
\end{table}
\FloatBarrier

\noindent
Test keys for Table~\ref{tab:nfr-trace-a}:
LF1 = NF.LF.1, LF2 = NF.LF.2, US1 = NF.US.1,
US2 = NF.US.2, PF1 = NF.PF.1, PF2 = NF.PF.2.

\begin{table}[H]
\centering
\small
\caption{Non-Functional Requirements Traceability Matrix Part 2}
\label{tab:nfr-trace-b}
\begin{tabularx}{\textwidth}{|c|*{6}{>{\centering\arraybackslash}X|}}
\hline
Req ID & PF3 & OP1 & OP2 & SEC1 & SEC2 & COM1 \\
\hline
SAF.R.1 & $\times$ & & & & $\times$ & \\
\hline
SAF.R.2 & $\times$ & & & & & \\
\hline
SOE.R.1 & $\times$ & & & & & \\
\hline
SOE.R.2 & $\times$ & & & & & \\
\hline
LON.R.1 & $\times$ & & & & & \\
\hline
LON.R.2 & $\times$ & & & & & \\
\hline
EPE.R.1 & & $\times$ & & & & \\
\hline
EPE.R.2 & & $\times$ & & & & \\
\hline
EPE.R.3 & & $\times$ & & & & \\
\hline
EPE.R.4 & & $\times$ & & & & \\
\hline
WER.R.1 & & $\times$ & & & & \\
\hline
WER.R.2 & & $\times$ & & & & \\
\hline
IAS.R.1 & & & $\times$ & & & \\
\hline
IAS.R.2 & & & $\times$ & & & \\
\hline
IAS.R.3 & & & $\times$ & & & \\
\hline
IAS.R.4 & & & $\times$ & & & \\
\hline
PRD.R.1 & & & $\times$ & & & \\
\hline
PRD.R.2 & & & $\times$ & & & \\
\hline
PRD.R.3 & & & $\times$ & & & \\
\hline
REL.R.1 & & & $\times$ & & & \\
\hline
REL.R.2 & & & $\times$ & & & \\
\hline
ACS.R.1 & & & & $\times$ & & \\
\hline
ACS.R.2 & & & & $\times$ & $\times$ & \\
\hline
INT.R.1 & & & & $\times$ & & \\
\hline
PRIV.R.1 & & & & $\times$ & & \\
\hline
IMM.R.1 & & & & $\times$ & & \\
\hline
IMM.R.2 & & & & $\times$ & & \\
\hline
LGL.R.1 & & & & & & $\times$ \\
\hline
LGL.R.2 & & & & & & $\times$ \\
\hline
STDCOMP.R.1 & & & & & & $\times$ \\
\hline
STDCOMP.R.2 & & & & & & $\times$ \\
\hline
\end{tabularx}
\end{table}
\FloatBarrier

\noindent
Test keys for Table~\ref{tab:nfr-trace-b}:
PF3 = NF.PF.3, OP1 = NF.OP.1, OP2 = NF.OP.2,
SEC1 = NF.SEC.1, SEC2 = NF.SEC.2, COM1 = NF.COM.1.

\section{Unit Test Description}

This section describes the planned unit tests for Proxi after the MIS
was completed. The goal of these tests is to check important module
logic before moving to integration and system testing.

Our general testing approach is simple. For each selected module, we
test the normal case first, then important edge cases and error cases.
Most of these tests will be automated. For modules that depend on
outside services or tools, mocks will be used instead of the real
service.

To keep the unit test section readable, each planned test is given a
short unit-test ID for traceability and a descriptive test name for
implementation.

\subsection{Unit Testing Scope}

The main modules in scope for direct unit testing are BH-Input,
BH-NLU, BH-Plan, BH-Safety, BH-Feedback, and BH-UI.

These modules were chosen because they contain important logic, are
easy to test in isolation, and have clear access programs in the
MIS.

The following are outside direct unit test scope: external speech
services, external MCP tools, browser, email, calendar, Notion, and
operating-system behaviour. These will be verified through mocks,
integration tests, and system tests instead.

Some support modules, such as SD-Types, SD-Store, SD-AIClient,
SD-ToolRegistry, and SD-Log, are lower priority for direct unit
testing in this document. They are either simple support modules or
are tested indirectly through higher-priority modules.

\subsection{Tests for Functional Requirements}

Most unit verification will be done using automated tests. The test
cases below are grouped by module and use short unit-test IDs so
they can be traced later in the report.

\subsubsection{Tests for BH-Input}

BH-Input is unit tested because it handles voice and text input,
state changes, and input retrieval.

\begin{enumerate}

\item \textbf{UT.INPUT.1 --
\texttt{test\_bh\_input\_init\_sets\_idle\_state}}

\textbf{Type:} Functional, Dynamic, Automatic

\textbf{Initial State:} Module not initialized.

\textbf{Input/Condition:} Call \texttt{initInput()}.

\textbf{Output/Result:} State becomes Idle, mode becomes TextOnly, and stored
text is empty.

\textbf{How test will be performed:} An automated unit test will call
\texttt{initInput()} and check the returned or stored state.

\item \textbf{UT.INPUT.2 --
\texttt{test\_bh\_input\_start\_capture\_from\_idle}}

\textbf{Type:} Functional, Dynamic, Automatic

\textbf{Initial State:} Module initialized and idle.

\textbf{Input/Condition:} Call \texttt{startCapture(Mixed)}.

\textbf{Output/Result:} State becomes Listening and mode becomes Mixed.

\textbf{How test will be performed:} An automated unit test will start capture
and check the new state.

\item \textbf{UT.INPUT.3 --
\texttt{test\_bh\_input\_duplicate\_start\_raises\_error}}

\textbf{Type:} Functional, Dynamic, Automatic

\textbf{Initial State:} Module already listening.

\textbf{Input/Condition:} Call \texttt{startCapture(VoiceOnly)} again.

\textbf{Output/Result:} Exception \texttt{AlreadyCapturing} is raised.

\textbf{How test will be performed:} An automated unit test will start capture
once and then call it again.

\item \textbf{UT.INPUT.4 --
\texttt{test\_bh\_input\_stop\_capture\_returns\_idle}}

\textbf{Type:} Functional, Dynamic, Automatic

\textbf{Initial State:} Module listening or processing.

\textbf{Input/Condition:} Call \texttt{stopCapture()}.

\textbf{Output/Result:} State returns to Idle.

\textbf{How test will be performed:} An automated unit test will start capture,
stop it, and verify the resulting state.

\item \textbf{UT.INPUT.5 --
\texttt{test\_bh\_input\_get\_last\_text\_raises\_when\_empty}}

\textbf{Type:} Functional, Dynamic, Automatic

\textbf{Initial State:} Module initialized with no stored text.

\textbf{Input/Condition:} Call \texttt{getLastText()}.

\textbf{Output/Result:} Exception \texttt{NoInputAvailable} is raised.

\textbf{How test will be performed:} An automated unit test will initialize the
module and request the last text right away.

\end{enumerate}

\subsubsection{Tests for BH-NLU}

BH-NLU is unit tested because it converts input text into intents and
parameters, and its behaviour is mostly deterministic.

\begin{enumerate}

\item \textbf{UT.NLU.1 --
\texttt{test\_bh\_nlu\_parse\_email\_command}}

\textbf{Type:} Functional, Dynamic, Automatic

\textbf{Initial State:} Parser available.

\textbf{Input/Condition:} Call \texttt{parseText("Draft an email to John about
tomorrow's meeting")}.

\textbf{Output/Result:} Returned intent matches draft-email and includes useful
parameters.

\textbf{How test will be performed:} An automated unit test will compare the
returned intent type and fields to the expected values.

\item \textbf{UT.NLU.2 --
\texttt{test\_bh\_nlu\_parse\_calendar\_command}}

\textbf{Type:} Functional, Dynamic, Automatic

\textbf{Initial State:} Parser available.

\textbf{Input/Condition:} Call \texttt{parseText("Create a calendar event for
Friday at 2 PM called Team Sync")}.

\textbf{Output/Result:} Returned intent matches create-calendar-event and
includes time and title information.

\textbf{How test will be performed:} An automated unit test will verify the
returned type and parsed fields.

\item \textbf{UT.NLU.3 --
\texttt{test\_bh\_nlu\_unsupported\_command\_raises\_parse\_error}}

\textbf{Type:} Functional, Dynamic, Automatic

\textbf{Initial State:} Parser available.

\textbf{Input/Condition:} Call \texttt{parseText("Paint my room blue")}.

\textbf{Output/Result:} Exception \texttt{ParseError} is raised.

\textbf{How test will be performed:} An automated unit test will give an
unsupported command and check the exception.

\end{enumerate}

\subsubsection{Tests for BH-Plan}

BH-Plan is unit tested because it creates plans from intents and
controls execution behaviour. External tools are mocked in these
tests.

\begin{enumerate}

\item \textbf{UT.PLAN.1 --
\texttt{test\_bh\_plan\_init\_sets\_default\_state}}

\textbf{Type:} Functional, Dynamic, Automatic

\textbf{Initial State:} Module not initialized.

\textbf{Input/Condition:} Call \texttt{initPlan()}.

\textbf{Output/Result:} Current plan is null and last status is Pending.

\textbf{How test will be performed:} An automated unit test will initialize the
module and check the stored values.

\item \textbf{UT.PLAN.2 --
\texttt{test\_bh\_plan\_create\_plan\_for\_supported\_intent}}

\textbf{Type:} Functional, Dynamic, Automatic

\textbf{Initial State:} Module initialized with a mocked valid tool registry.

\textbf{Input/Condition:} Call \texttt{planAction(i)} for a supported intent.

\textbf{Output/Result:} A plan is created with the expected tool and parameters.

\textbf{How test will be performed:} An automated unit test will mock the tool
registry and compare the returned plan.

\item \textbf{UT.PLAN.3 --
\texttt{test\_bh\_plan\_missing\_tool\_raises\_error}}

\textbf{Type:} Functional, Dynamic, Automatic

\textbf{Initial State:} Module initialized with no matching tool.

\textbf{Input/Condition:} Call \texttt{planAction(i)} for an unsupported tool
type.

\textbf{Output/Result:} Exception \texttt{NoToolFound} is raised.

\textbf{How test will be performed:} An automated unit test will use a mocked
registry with no valid match.

\item \textbf{UT.PLAN.4 --
\texttt{test\_bh\_plan\_execute\_plan\_updates\_status}}

\textbf{Type:} Functional, Dynamic, Automatic

\textbf{Initial State:} Module initialized with a valid mocked plan.

\textbf{Input/Condition:} Call \texttt{executePlan(p)}.

\textbf{Output/Result:} Status becomes Success or Failed based on the mocked
tool result.

\textbf{How test will be performed:} An automated unit test will mock tool
results and check the final status.

\item \textbf{UT.PLAN.5 --
\texttt{test\_bh\_plan\_cancel\_without\_pending\_plan\_fails}}

\textbf{Type:} Functional, Dynamic, Automatic

\textbf{Initial State:} Module initialized with no pending plan.

\textbf{Input/Condition:} Call \texttt{cancelPlan()}.

\textbf{Output/Result:} Exception \texttt{NoPendingPlan} is raised.

\textbf{How test will be performed:} An automated unit test will call
\texttt{cancelPlan()} after initialization and verify the exception.

\end{enumerate}

\subsubsection{Tests for BH-Safety}

BH-Safety is unit tested because it contains important decision rules
for risky actions and confirmations.

\begin{enumerate}

\item \textbf{UT.SAFE.1 --
\texttt{test\_bh\_safety\_low\_risk\_action\_is\_allowed}}

\textbf{Type:} Functional, Dynamic, Automatic

\textbf{Initial State:} Safety module available.

\textbf{Input/Condition:} Pass a low-risk action to the safety policy
functions.

\textbf{Output/Result:} The action is automatically allowed.

\textbf{How test will be performed:} An automated unit test will create a mock
low-risk action and verify the result.

\item \textbf{UT.SAFE.2 --
\texttt{test\_bh\_safety\_high\_irreversible\_action\_is\_blocked}}

\textbf{Type:} Functional, Dynamic, Automatic

\textbf{Initial State:} Safety module available.

\textbf{Input/Condition:} Pass a high-risk irreversible action to the safety
policy functions.

\textbf{Output/Result:} The action is blocked.

\textbf{How test will be performed:} An automated unit test will create a mock
high-risk action and verify the result.

\item \textbf{UT.SAFE.3 --
\texttt{test\_bh\_safety\_medium\_risk\_asks\_for\_confirmation}}

\textbf{Type:} Functional, Dynamic, Automatic

\textbf{Initial State:} Safety module available with mocked user response.

\textbf{Input/Condition:} Pass a medium-risk action to the confirmation path.

\textbf{Output/Result:} The action requires user confirmation and follows the
mocked yes or no response.

\textbf{How test will be performed:} An automated unit test will mock the user
response and verify the returned decision.

\item \textbf{UT.SAFE.4 --
\texttt{test\_bh\_safety\_confirmation\_timeout\_raises\_error}}

\textbf{Type:} Functional, Dynamic, Automatic

\textbf{Initial State:} Safety module available with no user response.

\textbf{Input/Condition:} Pass an action that needs confirmation to the
confirmation path.

\textbf{Output/Result:} The module raises the expected timeout error.

\textbf{How test will be performed:} An automated unit test will simulate no
response and verify the timeout behaviour.

\end{enumerate}

\subsubsection{Tests for BH-Feedback}

BH-Feedback is unit tested because it controls queue behaviour,
delivery state, and cancellation logic.

\begin{enumerate}

\item \textbf{UT.FEED.1 --
\texttt{test\_bh\_feedback\_init\_sets\_idle\_and\_empty\_queue}}

\textbf{Type:} Functional, Dynamic, Automatic

\textbf{Initial State:} Module not initialized.

\textbf{Input/Condition:} Call \texttt{initFeedback()}.

\textbf{Output/Result:} Queue is empty and status is Idle.

\textbf{How test will be performed:} An automated unit test will initialize the
module and check the queue and status.

\item \textbf{UT.FEED.2 --
\texttt{test\_bh\_feedback\_queue\_message\_appends\_successfully}}

\textbf{Type:} Functional, Dynamic, Automatic

\textbf{Initial State:} Module initialized.

\textbf{Input/Condition:} Call \texttt{queueMessage("Ready", Both)}.

\textbf{Output/Result:} The message is added to the queue.

\textbf{How test will be performed:} An automated unit test will queue one
message and check the queue size.

\item \textbf{UT.FEED.3 --
\texttt{test\_bh\_feedback\_speak\_now\_updates\_completed\_state}}

\textbf{Type:} Functional, Dynamic, Automatic

\textbf{Initial State:} Module initialized with mocked output channels.

\textbf{Input/Condition:} Call \texttt{speakNow("Done", Both)}.

\textbf{Output/Result:} The final status becomes Completed on success.

\textbf{How test will be performed:} An automated unit test will mock the
output channels and check the final status.

\item \textbf{UT.FEED.4 --
\texttt{test\_bh\_feedback\_tts\_failure\_is\_handled}}

\textbf{Type:} Functional, Dynamic, Automatic

\textbf{Initial State:} Module initialized with a failing mocked TTS path.

\textbf{Input/Condition:} Call \texttt{speakNow("Done", VoiceOnly)}.

\textbf{Output/Result:} The expected failure path or exception is triggered.

\textbf{How test will be performed:} An automated unit test will mock TTS
failure and verify the result.

\item \textbf{UT.FEED.5 --
\texttt{test\_bh\_feedback\_cancel\_clears\_queue}}

\textbf{Type:} Functional, Dynamic, Automatic

\textbf{Initial State:} Module initialized with queued messages.

\textbf{Input/Condition:} Call \texttt{cancelAll()}.

\textbf{Output/Result:} Queue is empty and status returns to Idle.

\textbf{How test will be performed:} An automated unit test will queue
messages, cancel them, and check the final state.

\end{enumerate}

\subsubsection{Tests for BH-UI}

BH-UI is unit tested for state changes and message handling. Visual
quality itself is checked mainly by system tests.

\begin{enumerate}

\item \textbf{UT.UI.1 --
\texttt{test\_bh\_ui\_init\_sets\_idle\_state}}

\textbf{Type:} Functional, Dynamic, Automatic

\textbf{Initial State:} UI module not initialized.

\textbf{Input/Condition:} Call \texttt{initUI()}.

\textbf{Output/Result:} State becomes Idle and last message is empty.

\textbf{How test will be performed:} An automated unit test will initialize the
module and check the state.

\item \textbf{UT.UI.2 --
\texttt{test\_bh\_ui\_update\_view\_changes\_ui\_state}}

\textbf{Type:} Functional, Dynamic, Automatic

\textbf{Initial State:} UI module initialized.

\textbf{Input/Condition:} Call \texttt{updateView(Listening)}.

\textbf{Output/Result:} UI state changes to Listening.

\textbf{How test will be performed:} An automated unit test will call
\texttt{updateView()} and verify the new state.

\item \textbf{UT.UI.3 --
\texttt{test\_bh\_ui\_show\_message\_stores\_message}}

\textbf{Type:} Functional, Dynamic, Automatic

\textbf{Initial State:} UI module initialized.

\textbf{Input/Condition:} Call \texttt{showMessage("Draft created")}.

\textbf{Output/Result:} The stored message is updated and the UI state becomes
Displaying.

\textbf{How test will be performed:} An automated unit test will call
\texttt{showMessage()} and check the resulting state.

\item \textbf{UT.UI.4 --
\texttt{test\_bh\_ui\_prompt\_user\_returns\_response}}

\textbf{Type:} Functional, Dynamic, Automatic

\textbf{Initial State:} UI module initialized with mocked user response.

\textbf{Input/Condition:} Call \texttt{promptUser("Send this email?")}.

\textbf{Output/Result:} The prompt path returns the expected user response.

\textbf{How test will be performed:} An automated unit test will mock the user
response and verify the returned value.

\item \textbf{UT.UI.5 --
\texttt{test\_bh\_ui\_clear\_screen\_resets\_display}}

\textbf{Type:} Functional, Dynamic, Automatic

\textbf{Initial State:} UI module initialized with a shown message.

\textbf{Input/Condition:} Call \texttt{clearScreen()}.

\textbf{Output/Result:} State returns to Idle and the message is cleared.

\textbf{How test will be performed:} An automated unit test will show a
message, clear the screen, and verify the final state.

\end{enumerate}

\subsection{Tests for Non-Functional Requirements}

Most non-functional requirements are checked better by integration and
system testing. Only a few simple unit-level checks are included here.

\subsubsection{Tests for BH-Feedback}

\begin{enumerate}

\item \textbf{UT.NF.FEED.1 --
\texttt{test\_bh\_feedback\_cancel\_leaves\_no\_queued\_messages}}

\textbf{Type:} Non-Functional, Dynamic, Automatic

\textbf{Initial State:} Feedback module initialized.

\textbf{Input/Condition:} Queue a few messages and then call
\texttt{cancelAll()}.

\textbf{Output/Result:} Queue becomes empty and status returns to Idle.

\textbf{How test will be performed:} An automated unit test will queue
messages, cancel them, and check the final state.

\end{enumerate}

\subsubsection{Tests for BH-Safety}

\begin{enumerate}

\item \textbf{UT.NF.SAFE.1 --
\texttt{test\_bh\_safety\_high\_risk\_action\_is\_never\_auto\_approved}}

\textbf{Type:} Non-Functional, Dynamic, Automatic

\textbf{Initial State:} Safety module available.

\textbf{Input/Condition:} Pass a high-risk irreversible action to the safety
policy functions.

\textbf{Output/Result:} The action is blocked or requires confirmation, but it
is never auto-approved.

\textbf{How test will be performed:} An automated unit test will create a mock
high-risk action and verify the result.

\end{enumerate}

\FloatBarrier

\bibliographystyle{plainnat}

\bibliography{../../refs/References}

\newpage

\section{Appendix}

This section includes supporting information used by multiple V\&V
activities.

\subsection{Symbolic Parameters}

The following symbolic parameters centralize repeated thresholds and test-set
sizes used throughout this V\&V plan. Using these names instead of repeating
raw numeric values improves consistency, traceability, and maintainability.

\begin{table}[H]
\centering
\caption{Symbolic Parameters for V\&V Activities}
\label{tab:symbolic-params}
\begin{tabularx}{\textwidth}{|p{5.2cm}|p{2.7cm}|X|}
\hline
\textbf{Name} & \textbf{Value} & \textbf{Meaning} \\
\hline
\texttt{CMD\_SET\_SIZE} & 20 &
Number of representative commands in the main command acceptance set. \\
\hline
\texttt{CMD\_ACCEPT\_MIN} & 95\% &
Minimum acceptance rate for valid command input tests. \\
\hline
\texttt{STT\_CMD\_COUNT} & 10 &
Number of spoken command phrases in the STT dataset. \\
\hline
\texttt{STT\_SPEAKER\_COUNT} & 5 &
Number of speakers used in the STT sample set. \\
\hline
\texttt{STT\_ACC\_MIN} & 90\% &
Minimum acceptable speech-to-text word accuracy. \\
\hline
\texttt{INTENT\_SAMPLE\_COUNT} & 100 &
Number of labeled inputs used in intent recognition testing. \\
\hline
\texttt{INTENT\_ACC\_MIN} & 90\% &
Minimum acceptable intent recognition accuracy. \\
\hline
\texttt{TASK\_SUCCESS\_MIN} & 85\% &
Minimum success rate for representative end-to-end task execution. \\
\hline
\texttt{FEEDBACK\_DELAY\_MAX} & 2 s &
Maximum delay before visible or spoken feedback appears after completion
or known failure. \\
\hline
\texttt{FOLLOWUP\_SUCCESS\_MIN} & 90\% &
Minimum correct-resolution rate for conversational follow-up sequences. \\
\hline
\texttt{USABILITY\_PARTICIPANTS} & 10 &
Target number of participants in the main usability study. \\
\hline
\texttt{HIGHLEVEL\_TASK\_MIN} & 80\% &
Minimum successful completion rate for high-level tasks without low-level
system details. \\
\hline
\texttt{ORIENTATION\_TIME} & 5 min &
Maximum orientation period given before usability testing starts. \\
\hline
\texttt{TEXT\_RESIZE\_REQ} & 200\% &
Required text resize level during accessibility inspection. \\
\hline
\texttt{SESSION\_DURATION\_MIN} & 4 h &
Minimum continuous operation duration during robustness testing. \\
\hline
\texttt{MIC\_DISTANCE\_MAX} & 1 m &
Maximum expected user distance from the microphone in environment tests. \\
\hline
\texttt{ROOM\_TEMP\_MIN} & 15$^\circ$C &
Minimum room temperature in the expected physical environment test. \\
\hline
\texttt{ROOM\_TEMP\_MAX} & 30$^\circ$C &
Maximum room temperature in the expected physical environment test. \\
\hline
\texttt{SURVEY\_CLEAR\_MIN} & 4.0 / 5.0 &
Minimum acceptable mean score for survey items on clarity,
learnability, and satisfaction. \\
\hline
\texttt{SURVEY\_DAILY\_USE\_MIN} & 60\% &
Minimum proportion of participants who say they would consider daily use. \\
\hline
\end{tabularx}
\end{table}

\subsection{Usability Survey and Acceptance Criteria}

This survey is used as the primary post-session instrument for
evaluating perceived usability quality after task-based testing. It is
designed to validate the SRS goals for clarity, learnability,
responsiveness, accessibility, trust, and acceptance for daily use.

\paragraph{Administration Protocol.}
The survey is completed immediately after each participant finishes the
usability session. Participants are instructed to base ratings on their
own observed experience in that session only. Unless otherwise stated,
the quantitative questions use a 1--5 Likert scale where 1 indicates the
least favorable response and 5 indicates the most favorable response.

\vspace{0.5em}
\begin{enumerate}

\item \textbf{How easy was it to complete tasks using voice commands?}\\
Response Type: 1--5 Likert (Very Difficult $\rightarrow$ Very Easy).\\
Intent: Assess overall learnability and ease of first-use interaction.

\item \textbf{Were the system responses clear and understandable?}\\
Response Type: 1--5 Likert (Very Unclear $\rightarrow$ Very Clear).\\
Intent: Measure clarity of feedback, wording quality, and output usefulness.

\item \textbf{Did the system react quickly enough to your commands?}\\
Response Type: 1--5 Likert (Very Slow $\rightarrow$ Very Fast).\\
Intent: Capture user-perceived responsiveness and latency satisfaction.

\item \textbf{How intuitive did the interface feel while completing tasks?}\\
Response Type: 1--5 Likert (Very Confusing $\rightarrow$ Very Intuitive).\\
Intent: Evaluate navigation flow, control predictability, and consistency.

\item \textbf{Was visual and audio feedback easy to notice and understand?}\\
Response Type: 1--5 Likert (Not Accessible $\rightarrow$ Highly Accessible).\\
Intent: Evaluate multimodal feedback accessibility and comprehensibility.

\item \textbf{Did you feel confident that Proxi understood your intent?}\\
Response Type: 1--5 Likert (Never $\rightarrow$ Always).\\
Intent: Measure trust in recognition and intent interpretation behaviour.

\item \textbf{Did confirmations and safety prompts appear at the right time?}\\
Response Type: 1--5 Likert (Never Appropriate $\rightarrow$ Always\\
Appropriate).\\
Intent: Evaluate whether high-risk actions are protected without harming flow.

\item \textbf{How satisfied are you with the overall experience?}\\
Response Type: 1--5 Likert (Very Unsatisfied $\rightarrow$ Very Satisfied).\\
Intent: Measure end-to-end satisfaction and readiness for practical use.

\item \textbf{Would you consider using Proxi for daily tasks?}\\
Response Type: Yes / No.\\
Intent: Estimate adoption potential for routine personal productivity use.

\item \textbf{Which task or moment felt most confusing or frustrating?}\\
Response Type: Open text.\\
Intent: Identify concrete workflow breakdowns and their triggers.

\item \textbf{What one change would most improve your experience?}\\
Response Type: Open text.\\
Intent: Collect high-impact, actionable improvement directions.

\end{enumerate}

\paragraph{Interpretation and Acceptance Rules.}
Survey responses are analyzed at both item level and theme level.

\begin{itemize}
  \item Questions 1, 2, 4, 5, 6, and 8 must each achieve a mean score of at
  least \texttt{SURVEY\_CLEAR\_MIN}.
  \item Question 3 must also meet \texttt{SURVEY\_CLEAR\_MIN} and be
  compared against measured timing data from performance tests. If
  measured latency passes but perception is low, the issue is treated as
  a feedback and expectation-management problem; if both fail, it is
  treated as a performance defect.
  \item Question 7 is interpreted as a safety-usability balance check.
  Scores below \texttt{SURVEY\_CLEAR\_MIN} trigger targeted review of
  confirmation frequency, timing, and wording.
  \item For Question 9, at least \texttt{SURVEY\_DAILY\_USE\_MIN} of
  participants must answer Yes.
  \item Open-text responses (Questions 10 and 11) are coded into themes
  such as confusion, latency, recognition error, accessibility barrier,
  safety prompt friction, and missing guidance, then mapped to affected
  workflows.
  \item Any failed numeric threshold or repeated negative theme results
  in a tracked issue, assigned corrective action, and re-test of the
  revised workflow before final delivery when feasible.
\end{itemize}

\newpage
\section*{Appendix - Reflection}

\wss{This section is not required for CAS 741}

The information in this section will be used to evaluate the team members on
the graduate attribute of Lifelong Learning.

\input{../Reflection.tex}

\begin{enumerate}
  \item What went well while writing this deliverable?
  \item What pain points did you experience during this deliverable, and how
    did you resolve them?
  \item What knowledge and skills will the team collectively need to acquire
  to successfully complete the verification and validation of your project?
  Examples of possible knowledge and skills include dynamic testing
  knowledge, static testing knowledge, specific tool usage, Valgrind etc.
  You should look to identify at least one item for each team member.
  \item For each of the knowledge areas and skills identified in the
  previous question, what are at least two approaches to acquiring the
  knowledge or mastering the skill? Of the identified approaches, which
  will each team member pursue, and why did they make this choice?
\end{enumerate}

\newpage
\section*{Reflection - Team}

\textbf{1. What went well while writing this deliverable?}\\
The strongest part of this deliverable was that each person was able to work
from the same testing story even while focusing on different layers of the
project. Aman said the planning and validation sections felt most meaningful
when they were tied directly back to the SRS, because that kept the document
focused on whether Proxi would actually help a user rather than just whether
it was technically impressive. Ajay said the non-functional testing work was
valuable because it forced the team to define what ``works well'' should mean
for responsiveness, accessibility, and user experience, not just basic
correctness. Gourob noted that the functional traceability work became much
easier once each requirement had a clearly corresponding test case, which
gave the plan a coherent structure. Savinay pointed out that our coordination
improved compared with earlier deliverables because responsibilities were
communicated earlier, so dependent sections could be written with fewer
rewrites. Taken together, the deliverable went well because the team moved
beyond isolated writing and started treating the document like a single
engineering argument.

\vspace{0.75em}

\noindent\textbf{2. What pain points did you experience during this
deliverable, and how did you resolve them?}\\
The main challenge was balancing ambition with clarity. Aman described the
struggle as finding the point where the writing was specific enough to be
credible but not so dense that it became hard to follow. Ajay had a similar
issue on the non-functional side, where broad ideas such as usability or
accessibility had to be translated into testable conditions before they were
useful. Gourob emphasized the practical disruption caused by the late
submission period and the MSAF process, which created pressure and made the
team reorganize quickly once relief had been granted. Savinay added that the
team had to recover from a period where several members were unavailable near
the deadline, and that uncertainty made planning more difficult than the
technical writing itself. The common resolution was a tighter review process:
we clarified wording, narrowed vague claims into measurable test cases, and
reassigned remaining work so the document could still read as one consistent
plan. In practice, that meant using review passes to reduce noise, improve
consistency, and keep the final text grounded in evidence.

\vspace{0.75em}

\noindent\textbf{3. What knowledge and skills will the team collectively
need to acquire to successfully complete the verification and validation of
your project?}\\
The team needs a broader mix of testing and evaluation skills than a typical
single-layer software project. Aman identified automated testing,
coverage-oriented thinking, survey analysis, accessibility evaluation, and
evidence organization as especially important, because Proxi will only be
convincing if its claims are backed by clear measurements. Ajay highlighted
dynamic testing, static testing, automation tools, UI validation, and
accessibility-focused inspection, since the interface has to stay usable as
the project evolves. Gourob focused on AI-supported evaluation, especially
intent recognition and speech input, because those behaviours are core to the
system and cannot be judged by intuition alone. Savinay emphasized test
planning, scenario design, user-group selection, and the ability to decide
whether evidence is strong enough to support a claim, since Proxi depends on
several interacting components and requires more than one kind of validation.
Across the team, the shared gap is not just technical testing knowledge but
the ability to connect observations to engineering claims in a disciplined
way.

\vspace{0.75em}

\noindent\textbf{4. For each of the knowledge areas and skills identified in
the previous question, what are at least two approaches to acquiring the
knowledge or mastering the skill? Of the identified approaches, which will
each team member pursue, and why did they make this choice?}\\
The team sees two main ways to build these skills. The first is structured
study: course notes, lab material, documentation, and examples from prior
testing assignments. The second is practice on the actual project: building
test scripts, running trial validation sessions, checking outputs, and
adjusting the plan based on what the evidence shows. Aman said he learns
best when theory is connected to something concrete, so he will use course
material as a reference but focus on building and reviewing real validation
artifacts. Ajay preferred small prototype checks and direct UI testing,
because practical runs make the tooling and interface tradeoffs easier to
understand than reading alone. Gourob planned to rely more heavily on
implementation and analysis of real outputs, since he already has a strong
technical base and learns best by measuring the system directly. Savinay said
he would use both study and examples from open-source projects, but expects
the applied examples to be most useful because they show how tests, evidence,
and quality checks fit together in practice. The team’s overall choice is to
combine both approaches, using theory to set the standard and practice to
prove that the standard is met.

\end{document}